\section{Theoretical Proofs and Extensions}
\label{sec:theory_appendix}

\subsection{Supporting Propositions and Lemmas}

\begin{proposition}[Rank-Normalized Stacking Invariance]
\label{prop:rank_invariance}
Let detector $m$ produce scores $s_m(x)$, and let $h_m:\mathbb{R}\to\mathbb{R}$ be any strictly increasing transformation. Define the normalized rank feature
\[
r_m(x)=\frac{1}{n}\sum_{i=1}^n \mathbf{1}\{s_m(x_i)\leq s_m(x)\}.
\]
Then replacing $s_m(x)$ with $h_m(s_m(x))$ leaves $r_m(x)$ unchanged for every sample $x$, up to tie-handling. Therefore, any ELARA stacker using only detector-wise rank-normalized features is invariant to monotone score rescaling, calibration drift, or sensor-scale drift.
\end{proposition}
\begin{proof}
For any two samples $(x_i,x_j)$,
\[
s_m(x_i)\leq s_m(x_j)
\]
if and only if
\[
h_m(s_m(x_i))\leq h_m(s_m(x_j)),
\]
because $h_m$ is strictly increasing. Therefore, the ordering of all samples under detector $m$ is unchanged. Since $r_m(x)$ depends only on the ordering of scores and not on their numerical scale, $r_m(x)$ is unchanged. Any downstream stacker that receives only $r_m(x)$ is therefore invariant to such monotone transformations.
\end{proof}

\medskip
\noindent\textbf{Interpretation.}
This proposition justifies the rank-normalized stacking design. It protects the router from meaningless score-scale changes while preserving anomaly ranking information.

\begin{proposition}[Degenerate-Channel Guard for Saturation]
\label{prop:degenerate_guard}
Let a fused anomaly score be
\[
g(x)=\sum_{m=1}^M w_m s_m(x),
\qquad w_m\geq 0,\qquad \sum_m w_m=1.
\]
If detector $k$ is saturated on validation data so that
\[
\mathrm{Var}_V(s_k)\leq \epsilon,
\]
then replacing $w_k$ by $0$ and renormalizing the remaining weights changes pairwise score differences by at most $O(\sqrt{\epsilon})$. In the limit $\epsilon\to 0$, the saturated detector contributes no ranking information, and removing it cannot remove meaningful rank signal.
\end{proposition}
\begin{proof}
For any two samples $(x_i,x_j)$,
\[
|s_k(x_i)-s_k(x_j)|
\leq
O(\sqrt{\epsilon})
\]
by the bounded-variance assumption. The contribution of detector $k$ to the fused pairwise difference is
\[
w_k(s_k(x_i)-s_k(x_j)),
\]
which is also $O(\sqrt{\epsilon})$. Therefore, as $\epsilon\to 0$, detector $k$'s contribution to every pairwise ranking comparison vanishes. Removing $k$ and renormalizing the remaining non-degenerate detectors does not remove meaningful rank information from the fused score.
\end{proof}

\medskip
\noindent\textbf{Interpretation.}
This supports the degenerate-channel guard for saturated detectors. A detector that produces nearly constant scores should not influence routing or fusion.

\begin{proposition}[Inverted Detector Harm Under Positive Fusion]
\label{prop:inverted_harm}
Assume the clean fused margin between an anomalous and normal sample is
\[
D_0=g_0(X^+)-g_0(X^-),
\]
with $\mathbb{E}[D_0]>0$. Suppose an inverted detector $k$ has margin
\[
D_k=s_k(X^+)-s_k(X^-)
\]
with $\mathbb{E}[D_k]<0$. For a positive fusion weight $w>0$, the fused margin
\[
D_w=(1-w)D_0+wD_k
\]
satisfies
\[
\mathbb{E}[D_w]
=
(1-w)\mathbb{E}[D_0]+w\mathbb{E}[D_k]
<
\mathbb{E}[D_0].
\]
Under a Gaussian equal-variance margin model, AUROC is monotone increasing in the standardized margin mean. Therefore, assigning positive weight to an inverted detector decreases AUROC relative to zeroing it.
\end{proposition}
\begin{proof}
The expectation follows from linearity:
\[
\mathbb{E}[D_w]
=
(1-w)\mathbb{E}[D_0]+w\mathbb{E}[D_k].
\]
Since $\mathbb{E}[D_k]<0<\mathbb{E}[D_0]$, we have $\mathbb{E}[D_w]<\mathbb{E}[D_0]$.
Under the Gaussian equal-variance margin model,
\[
D_w\sim \mathcal{N}(\mu_w,\sigma^2),
\]
and
\[
\mathrm{AUROC}(g_w)=\Pr[D_w>0]=\Phi\left(\frac{\mu_w}{\sigma}\right).
\]
Because $\Phi$ is strictly increasing and $\mu_w<\mu_0$, positive weight on the inverted detector decreases AUROC.
\end{proof}

\medskip
\noindent\textbf{Interpretation.}
This proposition supports validation-based inversion guards. If a detector is reliably below chance, positive-weight fusion can harm ranking.

\begin{proposition}[Meta-Router Generalization Requires Enough Tasks]
\label{prop:generalization_tasks}
Let $\mathcal{H}$ be a finite class of meta-routers, and let the task-level loss $\ell(h,\tau)\in[0,1]$ denote normalized rank regret or negative-transfer loss. Given $T$ independent meta-training tasks, with probability at least $1-\delta$, every router $h\in\mathcal{H}$ satisfies
\[
L(h)
\leq
\widehat{L}_T(h)
+
\sqrt{\frac{\log|\mathcal{H}|+\log(1/\delta)}{2T}},
\]
where $L(h)=\mathbb{E}_\tau[\ell(h,\tau)]$ is the expected held-out task loss and $\widehat{L}_T(h)=\frac{1}{T}\sum_{t=1}^T\ell(h,\tau_t)$ is the empirical meta-training loss.
\end{proposition}
\begin{proof}
For any fixed $h$, Hoeffding's inequality gives
\[
\Pr\left(
L(h)-\widehat{L}_T(h)>\epsilon
\right)
\leq
\exp(-2T\epsilon^2).
\]
Applying a union bound over all $h\in\mathcal{H}$,
\[
\Pr\left(
\exists h\in\mathcal{H}: L(h)-\widehat{L}_T(h)>\epsilon
\right)
\leq
|\mathcal{H}|\exp(-2T\epsilon^2).
\]
Set the right-hand side equal to $\delta$:
\[
|\mathcal{H}|\exp(-2T\epsilon^2)=\delta.
\]
Solving for $\epsilon$ gives
\[
\epsilon=
\sqrt{\frac{\log|\mathcal{H}|+\log(1/\delta)}{2T}}.
\]
Thus, with probability at least $1-\delta$, the stated bound holds for every $h\in\mathcal{H}$.
\end{proof}

\medskip
\noindent\textbf{Interpretation.}
This explains why ELARA needs many tasks and leave-family-out evaluation. A high-capacity router trained on too few tasks can overfit meta-level structure and fail to generalize.

\subsection{Old ELARA/RGA Source-System Theory Stack}

Below are the foundation certificates and bounds inherited from the ELARA/RGA source-system framework.

\begin{certificate}[Switching Certificate]
\label{cert:switch}
Let $\Delta_{\mathrm{fired}} = A(\mathrm{route}) - A(\mathrm{default})$ be the routed-minus-default performance on the subset of inputs where the router fires. The router may replace the default \emph{only if} the one-sided level-$\alpha$ lower confidence bound satisfies $\mathrm{LCB}_{1-\alpha}(\Delta_{\mathrm{fired}}) > 0$. Under this rule the expected switch is non-negative at level $\alpha$: routing is admitted only with certified fired-set benefit, never on a point estimate.
\end{certificate}

\begin{diagnostic}[Mean-Reliability Dilution]
\label{diag:dilution}
Let per-expert reliabilities be $r_1,\dots,r_M$ with mean $\bar r=\tfrac1M\sum_m r_m$ and gate threshold $\tau$. There exist configurations with $\bar r > \tau$ while a minority expert collapses ($r_k\to 0$): mean-gating does not fire and the collapsed expert is still trusted. The miss probability is $P(\mathrm{miss})=\Phi\!\big((\bar\mu-\tau)/\bar\sigma\big)$; detecting minority collapse requires per-expert / min-gating, not the mean (Theorem~T3).
\end{diagnostic}

\begin{slack}[Router Generalisation Slack]
\label{prop:slack}
Let the router class have complexity $C$ (e.g.\ VC dimension or Rademacher complexity) and let the meta-training set contain $T$ tasks. The gap between meta-training and held-out expected rank-gain is bounded by $O\!\big(\sqrt{C/T}\big)$ with high probability. Hence the evidential strength of a routing claim requires $T \gtrsim C$: a high-capacity router over few tasks yields weak evidence even at a positive point estimate. This motivates a large, heterogeneous benchmark suite (the Gate U minimum-viable families) before any claim.
\end{slack}



